
\chapter{MATLAB Code Analysis: \texttt{complex2sphere.m}}
\author{}
\date{}


\maketitle

\section{Initial Overview}

As an AI, I am happy to help you understand this MATLAB script. I have analyzed your code, and here is a concise overview of its functionality based solely on the provided logic.

\paragraph{What is the main purpose of this file?}
The primary purpose of this file is to map coordinates from the 2D complex plane onto a 3D unit sphere. 

\paragraph{What type of analysis or calculation does it perform?}
It performs a geometric transformation using rational algebraic calculations. Specifically, it applies the mathematical formulas for a stereographic projection to convert Cartesian representations of complex numbers into 3D Cartesian coordinates $(x, y, z)$.

\paragraph{What is the mathematical/theoretical context?}
This code operates within the context of \textbf{complex analysis} and \textbf{geometry}, specifically representing the \textbf{Riemann sphere} (the extended complex plane). It establishes a bijection between the extended complex plane $\mathbb{C} \cup \{\infty\}$ and the unit sphere $S^2 \subset \mathbb{R}^3$.

\section{Step-by-Step Code Analysis}

\subsection*{1. Section/Function Title: Variable Initialization and Component Extraction}
\textbf{Explanation:} This section defines the function and extracts the real and imaginary components from the input complex number array $w$. 
The input $w$ can be a single complex number or an array of complex numbers $w = X + iY$. The \texttt{real()} and \texttt{imag()} functions isolate the $X$ and $Y$ coordinates on the complex plane, which are required for the 3D mapping calculations.

\begin{lstlisting}
function v=complex2sphere(w)
X=real(w); Y=imag(w);
\end{lstlisting}

\subsection*{2. Section/Function Title: Infinity/NaN Detection}
\textbf{Explanation:} In the context of the Riemann sphere, the concept of ``infinity'' is treated as a specific point. MATLAB often represents undefined or infinite complex operations as \texttt{NaN} (Not a Number). This step finds the linear indices of all elements in the input array $w$ that evaluate to \texttt{NaN}, saving them in the variable $I$ so they can be mapped to the ``point at infinity'' later.

\begin{lstlisting}
I=find(isnan(w));
\end{lstlisting}

\subsection*{3. Section/Function Title: Stereographic Projection Math}
\textbf{Explanation:} This block applies the core stereographic projection equations. It first calculates a shared denominator $R = 1 + X^2 + Y^2$. Then, it calculates the $x$, $y$, and $z$ components for the 3D unit sphere. 
Mathematically, the origin of the complex plane ($0 + 0i$) will map to the South Pole $(0, 0, -1)$, and the unit circle ($X^2 + Y^2 = 1$) will map to the equator ($z = 0$).

\begin{lstlisting}
R=1+X.^2+Y.^2;
x=2*X./R;
y=2*Y./R;
z=(-1+X.^2+Y.^2)./R;
\end{lstlisting}

\subsection*{4. Section/Function Title: Exception Handling and Output Construction}
\textbf{Explanation:} This final step resolves the \texttt{NaN} values detected earlier. It takes the indices stored in $I$ and forces their corresponding $x$, $y$, and $z$ coordinates to $0$, $0$, and $1$, respectively. This properly maps complex infinity to the North Pole of the Riemann sphere. Finally, it concatenates the $x$, $y$, and $z$ column vectors into a single output matrix $v$, representing the 3D coordinates.

\begin{lstlisting}
x(I)=0;y(I)=0;z(I)=1;
v=[x y z];
\end{lstlisting}

\section{Expected Outputs and Visuals}

This specific code is entirely computational; it does \textbf{not} inherently generate, plot, or output any visual graphs to the MATLAB figure window, nor does it log formatted outputs to the console (aside from the standard matrix return if the semicolon is omitted during the function call). It returns an $N \times 3$ matrix of numbers representing spatial coordinates.

However, if you were to visualize the output matrix $v$ using MATLAB plotting functions (such as \texttt{scatter3} or \texttt{plot3}):
\begin{itemize}
    \item \textbf{Axes:} The axes would be the 3D Cartesian coordinates $X$, $Y$, and $Z$ bounded roughly between $-1$ and $1$.
    \item \textbf{Visual Representation:} The plot would display points perfectly distributed along the surface of a 3D sphere with a radius of $1$, centered at the origin $(0,0,0)$. 
    \item \textbf{Mathematical Translation:} Inputs near the origin ($0$) in the complex plane would visually cluster near the bottom of the sphere (South Pole, $z=-1$). Inputs with large absolute values in the complex plane would map closer to the top of the sphere (North Pole, $z=1$). Inputs explicitly passed as \texttt{NaN} would be plotted exactly at the apex $(0,0,1)$.
\end{itemize}
